\textbf{二阶逻辑}是一阶逻辑的扩展，它引入了可被量化的关系变元和函数变元。
二阶逻辑的结构同一阶结构一样，给出语言中所有非逻辑符号的解释。
然而，变元指派还要将论域上的关系和函数指派给二阶变元。
二阶公式的满足关系的定义方式与一阶情形相同，但扩充为可处理二阶变元和量词。

二阶量词使得二阶逻辑的\textbf{表达能力}强于一阶逻辑。
例如，可以不使用~$\eq$，而通过 $\lforall[X][(X(x) \liff X(y))]$ 来定义结构论域上的等词关系。
二阶逻辑可以表达关系的\textbf{传递闭包}，而这在一阶逻辑中是不可表达的。
二阶量词还能表达论域的性质，即论域是有穷的还是无穷的，是可数的还是不可数的。
例如，这意味着存在一个二阶语句~$\fn{Inf}$，使得 $\Sat{M}{\fn{Inf}}$ 当且仅当 $\Domain{M}$ 是无穷的。
重要的是，这些是\textbf{纯}二阶语句，即它们不含任何非逻辑符号。
由于紧致性定理和Löwenheim--Skolem定理，不存在具有这些性质的一阶语句。
这也表明，\textbf{紧致性定理和Löwenheim--Skolem定理在二阶逻辑中不成立}。

二阶量化还使得我们可以用单个语句取代一阶模式。
例如，\textbf{二阶算术}~$\Th{PA^2}$ 由~$\Th{Q}$ 的公理加上单个\textbf{归纳公理}\[
\lforall[X][((X(\Obj 0) \land \lforall[x][(X(x) \lif X(x'))]) \lif
  \lforall[x][X(x)])].
\]构成
与一阶~$\Th{PA}$ 不同，$\Th{PA^2}$ 的所有二阶模型都同构于标准模型。
换言之，$\Th{PA^2}$ \textbf{没有非标准模型}。

由于二阶逻辑包含一阶逻辑，它是不可判定的。
一阶逻辑至少是可公理化的，即它有一个可靠且完备的推导系统。
二阶逻辑则不然，它是\textbf{不可公理化的}。
因此，二阶逻辑的有效式集合是高度不可计算的。
事实上，纯二阶逻辑可以表达诸如\textbf{连续统假设}等集合论断言，而这些断言独立于集合论。